\section{Introduction}

Area-minimizing hypercones arise as tangent cones at singularities of
area-minimizing hypersurfaces and play a central role in their regularity
theory \cite{HardtSimon85,Simon94,Szekelyhidi26,FiresterTsiamisWang26}.
Recent work of Chodosh and collaborators \cite{CMS26,CMSW25}
establishes generic regularity in ambient dimensions nine, ten, and eleven
for Plateau solutions and homology minimizers. Their arguments rely on
the analysis of minimizing cones and Jacobi fields. In particular, the
treatment of cylindrical singularities in \cite{CMSW25} uses the strict
stability and strict area minimality of minimizing quadratic cone factors.
We study three closely connected aspects of algebraic minimal hypercones:
their variational properties, the structure of their associated Laplacian
algebras, and their construction through polynomial realizations of
algebraic closure relations. Laplacian algebras provide a common framework
for these questions: their closure identities encode both the minimality
equation for a defining polynomial and the differential expressions
governing its subcalibrations.

The construction of algebraic minimal hypercones and their classification
in low degree have a long history. In 1967, Hsiang~\cite{Hsiang67}
classified the quadratic case and initiated the study of cubic minimal
hypercones. To formulate these problems algebraically, write
\(N=n+1\), \(n\geq2\), \(r=|z|\), and \(R=r^2\), and set
\begin{equation}\label{eq:L-definition}
 \cL(H)=|\nabla H|^2\Delta H
 -\frac12\left\langle\nabla|\nabla H|^2,\nabla H\right\rangle,
\end{equation}
for \(H\in C^2\). If \(F\) is a square-free homogeneous polynomial
whose real irreducible factors have nonsingular real zeros,
the regular part of \(\{F=0\}\) is minimal exactly when
\begin{equation}\label{eq:intro-minimality}
 \cL(F)=\lambda_FF,
\end{equation}
for a polynomial \(\lambda_F\); see \cref{prop:divisibility}.

Without additional structural assumptions, both classification and
explicit construction remain difficult. For example, the
Cartan--M\"unzner identities describe
isoparametric families whose minimal members yield algebraic hypercones
\cite{Muenzer80,Muenzer81}. For radial minimal cubics, where
\(\lambda_F\) is proportional to \(R\), Peng--Xiao~\cite{PengXiao93} and, independently,
Tkachev~\cite{TkachevClassification10,TkachevClifford10}
developed a systematic classification theory.

Once a minimal hypercone is given, a further question is whether it
minimizes area. The classical minimizing quadratic cones
\cite{BDGG69,Lawson72,Simoes73} provide a model for an algebraic
approach to this question. Building on earlier calibration and
subsolution arguments, De Philippis and Paolini
\cite{DePhilippisPaolini} gave a short proof of the area minimality
of Simons cones using subcalibrations. Suitable divergence inequalities
on the two sides replace the divergence-free condition and give direct
perimeter comparisons; see \cref{rem:subcalibration-history}.

Motivated by these quadratic constructions, we seek potentials
\(G=aF\), where \(a\) is a positive homogeneous multiplier on an
appropriate conical domain. Its normalized gradient satisfies
\begin{equation}\label{eq:div-normalized-gradient}
 X_G=\frac{\nabla G}{|\nabla G|},\qquad
 \diver X_G=\frac{\cL(G)}{|\nabla G|^3},
\end{equation}
where \(\nabla G\ne0\). On the regular zero set in this domain,
\(\nabla G=a\nabla F\), so \(X_G\) agrees
with the exterior normal to \(\{F<0\}\). The following definition
includes the removable singular sets needed in our applications.

\begin{definition}[Homogeneous subcalibration potential]
\label{def:homogeneous-subcalibration-potential}
Let \(F\) be a nonzero homogeneous polynomial on \(\R^N\).
A continuous function \(G\) on \(\R^N\) is a
\emph{homogeneous subcalibration potential} for \(F\) if there are
a closed conical set \(Z\) with \(\cH^{N-1}(Z)=0\) and a positive
homogeneous function \(a\in C^2(\R^N\setminus Z)\) such that
\(G=aF\) on \(\R^N\setminus Z\), \(G\) is homogeneous of degree
\(\ell>0\), and
\begin{equation}\label{eq:master-sign}
 \nabla G\ne0,\qquad G\cL(G)\ge0
 \quad\text{on }\R^N\setminus Z.
\end{equation}
\end{definition}

Since polynomial sublevel sets have locally finite perimeter,
\cref{cor:gradient-subcalibration} gives perimeter comparison
for \(\{F<0\}\) across \(Z\).

For the Simons cone, with \(x,y\in\R^m\), \(m\ge4\), the potential
in \cite[Section~2]{DePhilippisPaolini} is, up to a positive constant,
\[
 F=|x|^2-|y|^2,\qquad
 G=|x|^4-|y|^4=(|x|^2+|y|^2)F=RF.
\]
For Lawson cones in \(\R^k\times\R^h\), \(k,h\ge2\), put
\(U=|x|^2/(k-1)\), \(V=|y|^2/(h-1)\), and \(F=U-V\).
In the dimension ranges of
\cite[Lemma~4.1, equation~(4.7)]{DePhilippisMaggi14}, the quartic
subcalibration potential factors as
\begin{equation}\label{eq:intro-Lawson-multiplier}
 G=\frac{U^2-V^2}{4}=aF,\qquad a=\frac{U+V}{4}.
\end{equation}
Phase-dependent alternatives are discussed in \cite{Liu19} and
\cref{rem:phase-local-fields}.
In the two polynomial examples above, both \(F\) and \(a\) lie in
\(\R[|x|^2,|y|^2]=\R[R,F]\). This algebra is closed under
\(\Delta\) and the gradient pairing
\(\Gamma(u,v)=\langle\nabla u,\nabla v\rangle\), so the
subcalibration calculation stays in the same two quadratic variables.

More generally, a subalgebra \(\cA\subset\R[z_1,\ldots,z_N]\) is
a \emph{Laplacian algebra} if \(R\in\cA\) and
\(\Delta\cA\subset\cA\).
It is graded and finitely generated
\cite[Corollary~22 and Lemma~24]{MendesRadeschi}, and the product
rule gives closure under \(\Gamma\).
For homogeneous generators \(I=(I_1,\ldots,I_k)\), closure is
recorded by polynomial identities
\[
 \Delta I_i=b_i(I),\qquad
 \Gamma(I_i,I_j)=c_{ij}(I),
\]
which form the \emph{closure table}
\cite[Proposition~37(a)]{MendesRadeschi}.
For \(F\in\cA\) and a homogeneous multiplier \(a=\alpha\circ I\),
with \(\alpha\) smooth on an open set \(W\subset\R^k\), the chain
rule expresses \(\cL(F)\), \(\cL(aF)\), and \(|\nabla(aF)|^2\)
in the generator variables on \(I^{-1}(W)\).
These closure tables may therefore guide the construction and
recognition of minimal algebraic hypercones; see
\cref{sec:gradient-square-closure}.

We use the structural theory of Mendes--Radeschi
\cite{MendesRadeschi,MendesRadeschiQuadratic,MendesRadeschiMaximality}
to study cubic and quartic models and develop new polynomial realizations.
We begin with algebras generated by the radius squared and one
defining polynomial.

\subsection{Two generators and radial multipliers}

In this two-generator case, the Mendes--Radeschi characterization
\cite[Proposition~37(c)]{MendesRadeschi} leads to an isoparametric
family. We construct polynomial subcalibrations for all of its
nonflat minimizing cones.

\begin{maintheorem}[Polynomial subcalibrations for isoparametric hypercones]
\label{main:isoparametric}
Let $C\subset\R^N$ be a nonflat area-minimizing isoparametric hypercone,
and let $F$ be its centered Cartan--M\"unzner defining polynomial, as
in \eqref{eq:iso-F-z}. There is an explicit homogeneous polynomial
$a_F$, positive on $\R^N\setminus\{0\}$, such that $G=a_FF$ satisfies
\[
 \deg G\le10,\qquad \nabla G\ne0
 \quad\text{on }\R^N\setminus\{0\},
\]
and
\[
 \operatorname{sgn}(G)\diver X_G>0\qquad\text{off }C.
\]
The field $X_G$ is a subcalibration of $C$, and its divergence proves
strict area minimality.
\end{maintheorem}
Lawlor \cite{Lawlor91}, Ferus--Karcher \cite{FerusKarcher85} and
Wang \cite{Wang94} established the minimizing results used here;
see \cref{sec:isoparametric-reduction} for their precise scope.
Our explicit potentials have a dimension-independent degree bound
and recover strict area minimality. The complete list is in
\cref{thm:iso-polynomial-table}.
For harmonic defining polynomials, a radial criterion applies in
arbitrary degree.

\begin{maintheorem}[A radial criterion for harmonic polynomials]
\label{main:radial-criterion}
Let $F\not\equiv0$ be a homogeneous polynomial of degree $e\ge2$
on $\R^N$, $N\ge3$, and suppose
\[
 \Delta F=0,\qquad S=|\nabla F|^2,\qquad \Gamma(F,S)=2qF
\]
for a polynomial $q$. If, for a constant $\kappa\ge0$,
\[
 Rq\le\kappa S\quad\text{on }\R^N,\qquad
 N\ge2e+1+2\sqrt\kappa,
\]
then $G=r^{(N-1-2e)/2}F$ generates a two-sided subcalibration.
The cone $C=\{F=0\}$, oriented as the boundary of $\{F<0\}$,
is strictly area minimizing, and
\[
 Q_C(u)\ge\left[\frac{(N-1-2e)^2}{4}-\kappa\right]
       \int_Cr^{-2}u^2,\qquad u\in C_c^\infty(C\setminus\{0\}).
\]
\end{maintheorem}
The stability estimate follows from the multiplier--Jacobi identity
and the homogeneous ground-state bound
(\cref{thm:Jacobi-trace,thm:ground-state-gap}).
A finite $\kappa$ exists exactly when $\nabla F$ vanishes only at
the origin. Its optimal value is at least $e(e-1)$, with equality
precisely for normalized harmonic Cartan--M\"unzner polynomials
(\cref{prop:harmonic-global-comparison,prop:harmonic-sharp-comparison}).

\begin{remark}[Realizability and regular algebraic cones]
\label{rem:regular-algebraic-realizability}
This regularity requirement connects \cref{main:radial-criterion}
with Perdomo's question about non-isoparametric regular algebraic
minimal hypersurfaces of spheres
\cite[Section~1 and Main theorem]{PerdomoRegular2006}; his nonexistence
theorem concerns irreducible cubics in \(\Sph^3\).
We conjecture that, for \(N\ge3\), a nonconstant real homogeneous
irreducible polynomial \(F\) whose zero link is nonempty and minimal,
with \(\nabla F\ne0\) everywhere on that link, defines an
isoparametric hypersurface. In the harmonic case, we expect \(F\)
itself, up to scaling, to be a harmonic Cartan--M\"unzner polynomial.
\Cref{prop:regular-radial-weight-rigidity} proves this harmonic
conclusion when \(q=\lambda R^{e-2}\), an additional condition on
the minimality weight, distinct from choosing the radial multiplier
\(a=r^\sigma\) in \cref{main:radial-criterion}.
\end{remark}

\subsection{Gradient-square closure and quadratic compositions}

Adjoining $S=|\nabla F|^2$ closes the algebra in a distinguished
class. For $N\ge3$, odd degree at least three then forces
$\deg F=3$. If the critical set consists only of the origin, $F$
is a Cartan cubic; the other cubic members are described in
\cref{prop:cubic-gradient-classification}.
In even degree, we record the closure identities and a short
isoparametric consequence on the regular link, followed by a matrix
quartic example.

For three-generator degrees $(2,2,g)$, $g>2$, subtracting a polynomial
in the two block radii from the higher-degree generator makes it
separately harmonic and bihomogeneous. The motivating OT--FKM splitting
quartics, developed from \cite{CuiFKM25}, are quadratic compositions
of Clifford coordinates, as are Riedler's quartic pullbacks.
\Cref{sec:two-quadratic} develops their general orthogonal-block criterion.
These compositions have generator degrees $(2,4,4)$, distinguished
from strict gradient-square closure $(2,4,6)$.

\begin{maintheorem}[Subcalibrations from orthogonal quartic blocks]
\label{main:bihomogeneous}\label{main:quartic-table}
Let $E_0,E_1,E_2$ be nonzero homogeneous quartic polynomials on
$\R^N$, with $R=|x|^2$, satisfying
\[
 E_0+E_1+E_2=R^2,\qquad
 \Gamma(E_i,E_j)=16RE_i\delta_{ij},\qquad
 \Delta E_i=8r_iR.
\]
Assume $r_i>0$, $r_1,r_2\ge2$, and either
\[
 r_1+r_2\ge9,
 \qquad\text{or}\qquad
 r_1=r_2=4\ \text{ and }\ r_0\ge5.
\]
Set $F=(r_2-1)E_1-(r_1-1)E_2$ and
$\Sigma=\{E_1=E_2=0\}$. Then there is a homogeneous polynomial
potential $G=AF$ of degree at most ten, with $A>0$ off $\Sigma$
and $\operatorname{Crit}G=\Sigma$, such that
\[
 \operatorname{sgn}(G)\diver X_G>0\qquad\text{where }F\ne0.
\]
The exceptional set $\Sigma$ is removable for this field. The
oriented boundary of $\{F<0\}$ has support $\{F=0\}$ and is
strictly area minimizing and strictly stable.
\end{maintheorem}
Together with the instability results, the theorem gives the minimizing
classification for the splitting family and recovers Riedler's
Clifford--Simons range. The same table
also treats the octonionic Hopf pairing and the balanced $(3,7,7)$
realization. Their weights, potentials and ranges are collected in
\cref{tab:quartic-weight-realizations}; their different singular
strata distinguish the geometric realizations. In particular,
\cref{thm:orthq-infinite-families} constructs two further infinite
families in $\R^{8k}$ with an unused block of weight two; they are
area minimizing exactly for $k\ge3$. Quadratic Jordan and Gram
coordinates also give the degree-independent profile equation of
\cref{prop:three-quadratic-profile-criterion} and the harmonic
sextic family in \cref{thm:complex-gram-cycle-sextic}.

\subsection{Cubic structure and a common subcalibration}

The historical and algebraic background is recalled in
\cref{sec:cubic-algebraic-theory}. For harmonic minimal cubics, we
obtain a sharp bound on the smallest algebra.
For \(\cA=\R[I_1,\ldots,I_k]\), its \emph{spherical quotient dimension}
is \(\dim I(\Sph^{N-1})\), where \(I=(I_1,\ldots,I_k)\).

\begin{maintheorem}[Five generators for harmonic minimal cubics]
\label{main:cubic-generator-bound}
Let $F$ be a nonzero irreducible homogeneous harmonic cubic on its
essential Euclidean space, with minimal regular zero set. The smallest
Laplacian algebra $\cA_R(F)$ has at most five homogeneous generators.
For a radial cubic at most four suffice. For a nonradial cubic the
algebra is polynomial on five algebraically independent generators of
degrees $(2,2,3,3,4)$, with spherical quotient dimension four;
its ambient dimension is $8$, $14$, or $26$.
\end{maintheorem}
Here radial means that $\cL(F)=-qF$ has $q$ proportional to $R$.
The intrinsic closure recovers a Cartan polynomial and forces the
dimension list $8,14,26$. Explicit Cartan deformations realize all
three dimensions: the models in $\R^8$ and $\R^{14}$ are unstable,
whereas the model in $\R^{26}$ is strictly area minimizing.
For radial cubics, the following criterion is read from $\Delta S$.

\begin{maintheorem}[A cubic trace criterion]
\label{main:cubic-trace}
Let $F\not\equiv0$ be a homogeneous harmonic cubic on $\R^N$,
$N\ge8$, satisfying $\Gamma(F,S)=4RF$. If
\[
 \Delta S\ge\max\left\{16,\frac{2(40-N)}3\right\}R,
\]
then there exists a continuous homogeneous potential $G=AF$ of
degree six, with $A>0$ on $\{S>0\}$ and $G$ smooth there, whose
gradient is nonzero there and whose normalized gradient has the
strict divergence sign of $F$ off its zero set. The critical set
$\{S=0\}$ is removable. Hence $\{F<0\}$ is locally perimeter
minimizing, and its oriented boundary has support $\{F=0\}$ and
is strictly area minimizing.
\end{maintheorem}
The stronger form, \cref{thm:cubic-unified-sextic}, also covers the
low-dimensional Clifford cases outside this sufficient trace bound.
It yields the common degree-six potential $F\sqrt{RS}$ for rank-four
trace cubics, complexified rank-three Jordan norms and the $G_2$
trilinear cubic. The nonradial model in $\R^{26}$ uses the same
potential with its five-generator defect calculation.

For the Clifford family, let $(A_0,\ldots,A_q)$ be a symmetric
Clifford system on $\R^{2m}$, $q\ge1$, and put
\[
 F(y,z)=\sum_{i=0}^qz_i\langle A_i y,y\rangle.
\]

\begin{maintheorem}[Clifford cubic classification]
\label{main:clifford-classification}
The oriented boundary of $\{F<0\}$ is area minimizing exactly when
$m\ge4$, $q\ge3$, and $(m,q)\ne(4,3)$. In that range it is
strictly area minimizing, and, after replacing $F$ by $F/2$,
its Laplacian data satisfy \cref{thm:cubic-unified-sextic}.
For $(m,q)=(4,3)$ the cone
is unstable although every tangent cone at a nonzero point is area
minimizing.
\end{maintheorem}
The classification uses the common positivity criterion,
singular-stratum obstructions and an explicit negative variation for
the exceptional pair. The examples and their criteria are summarized
in \cref{tab:intro-family-summary}.

\begin{table}[htbp]
\centering
\small
\setlength{\tabcolsep}{4pt}
\renewcommand{\arraystretch}{1.2}
\caption{Selected families and the criteria used for their minimizing
members. The third column counts generators of the smallest Laplacian
algebra $\cA_R(F)$; $\le4$ is an upper bound. Explicit potentials
are constructed in the cited results.}
\label{tab:intro-family-summary}
\begin{tabular}{@{}>{\raggedright\arraybackslash}p{.24\linewidth}p{.07\linewidth}p{.09\linewidth}>{\raggedright\arraybackslash}p{.25\linewidth}>{\raggedright\arraybackslash}p{.25\linewidth}@{}}
\toprule
Family & $\deg F$ & Number & Criterion & Minimizing range \\
\midrule
Isoparametric & $g$ & $2$ & Two-generator classification; \cref{main:isoparametric}
 & Full list in \cref{thm:iso-polynomial-table} \\
Orthogonal quartic families & $4$ & $3$ & Three-weight criterion; \cref{main:quartic-table}
 & Splitting, Hopf, pullback and balanced ranges in \cref{tab:quartic-weight-realizations} \\
Rank-four trace & $3$ & $3$ & Cubic trace bound; \cref{main:cubic-trace}
 & $N=27$; \cref{cor:quaternionic-trace-cubic} \\
Complexified rank-three Jordan norms & $3$ & $4$ & Cubic trace bound; \cref{main:cubic-trace}
 & $N=18,30,54$; \cref{thm:exc-jordan-minimizing} \\
$G_2$ trilinear & $3$ & $4$ & Cubic trace bound; \cref{main:cubic-trace}
 & $N=21$; \cref{thm:exc-g2} \\
Clifford cubics & $3$ & $\le4$ & Full cubic criterion; \cref{thm:cubic-unified-sextic}
 & $m\ge4$, $q\ge3$, $(m,q)\ne(4,3)$ \\
Nonradial Cartan deformation & $3$ & $5$ & Five-generator positivity; \cref{thm:nonradial-minimizing}
 & $N=26$ \\
Octonionic products; complex Pfaffians & $d\ge3$ & $d+1$ & Spectral threshold; \cref{main:spectral-criterion}
 & Every $d\ge3$; \cref{main:spectral-families} \\
Complex determinants & $m\ge3$ & $m+1$ & Minor-norm inequalities; \cref{sec:complex-determinants}
 & Every $m\ge3$; \cref{main:complex-determinants} \\
\bottomrule
\end{tabular}
\end{table}
\FloatBarrier
\subsection{Pullbacks and arbitrary-degree spectral families}

Riedler's pullback construction and question
\cite[Theorem~5.2.5 and Question~5.3.3(2)]{RiedlerThesis25}
motivate the transfer theorem in \cref{thm:clifford-area-transfer}.
A homogeneous target subcalibration $H$ with
$H\Delta H\ge-M|\nabla H|^2$ transfers through a quadratic Clifford
map $P$ to $R^aH(P)$ in sufficiently large admissible dimensions,
under the singular-set hypotheses of that theorem.
The corresponding smallest algebra has an exact description.

\begin{maintheorem}[The smallest algebra under Clifford pullback]
\label{main:clifford-smallest-algebra}
Let $A_1,\ldots,A_m$ be symmetric endomorphisms of $\R^N$, $m\ge2$,
with $A_iA_j+A_jA_i=2\delta_{ij}\Id$, and set
\[
 P(x)=(\langle A_1x,x\rangle,\ldots,\langle A_mx,x\rangle),
 \qquad R=|x|^2,\quad q(y)=|y|^2,\quad Q=q(P).
\]
For every nonzero essential homogeneous polynomial $f$ of degree
$s\ge2$ on $\R^m$,
\[
 \cA_R(f\circ P)=\R[R,P^*\cA_q(f)].
\]
If $Q<R^2$ somewhere, this algebra is isomorphic to the polynomial
extension $\cA_q(f)[t]$, with target degrees doubled and $\deg t=2$.
In this case, its minimal number of homogeneous generators is
exactly one more than that of $\cA_q(f)$.
\end{maintheorem}
Iteration doubles the defining degree at each step and, under the
nondegeneracy hypothesis, adds exactly one minimal generator per step
(\cref{cor:iterated-clifford-transfer}). Cubic applications appear in
\cref{cor:common-cubic-pullback,thm:cartan-pullback}.

For arbitrary degree, nonnegative spectral variables give a second
uniform construction. Write $e_k$ for elementary symmetric polynomials,
with $e_0=1$.

\begin{maintheorem}[A spectral criterion and its octonionic and Pfaffian applications]
\label{main:spectral-criterion}\label{main:spectral-families}
Let $F,E_1,\ldots,E_d$ be real homogeneous polynomials on $\R^N$
of degrees $d,2,\ldots,2d$, where $d\ge2$, $F\not\equiv0$, and
$E_1=R$. Assume that at every point $E_k=e_k(r_1,\ldots,r_d)$ for
some nonnegative $r_i$, and that
\[
 \begin{gathered}
 \Delta F=0,\qquad S:=\Gamma(F,F)=E_{d-1},\\
 \Gamma(F,E_k)=2(d-k+1)E_{k-1}F,\\
 \Gamma(E_i,E_j)=4\sum_{\alpha=1}^d
       r_\alpha\,\partial_\alpha e_i(r)\partial_\alpha e_j(r).
 \end{gathered}
\]
If
\[
 \Delta S\ge\left(24-\frac{16}{d}\right)E_{d-2}
 \qquad\text{on }\{S>0\},
\]
there exists a homogeneous potential $G=AF$ of degree $2d-1$,
with $A>0$ and $\nabla G\ne0$ there, satisfying
$G\cL(G)>0$ wherever $F\ne0$. The hypotheses force $N\ge4d$.
The critical set $\{S=0\}$ is removable, and the normalized gradient
of $G$ gives a two-sided subcalibration. The phase and boundary
current are locally minimizing. The cone is strictly area minimizing,
and its regular part satisfies
\[
 Q_C(u)\ge\frac{(N-4d+1)^2}{4}\int_C |x|^{-2}u^2,
 \qquad u\in C_c^\infty(C_{\reg}\setminus\{0\}).
\]

In particular, for every $d\ge2$, each of the following harmonic degree-$d$ cones
is the support of a strictly area-minimizing integral boundary
and has strictly stable regular part:
\begin{enumerate}[label=\textup{(\roman*)}]
\item $\{\operatorname{Re}(x_1\cdots x_d)=0\}\subset\mathbb O^d$,
for every fixed parenthesization;
\item $\{\operatorname{Re}\operatorname{Pf} Z=0\}$ on the complex
skew-symmetric $2d\times2d$ matrices, with Euclidean norm
$\sum_{i<j}|z_{ij}|^2$.
\end{enumerate}
Their critical sets have real codimension sixteen and twelve, respectively.
The Hardy lower bounds are $(4d+1)^2/4$ and
$(4d^2-6d+1)^2/4$, respectively.
\end{maintheorem}
The proof uses Newton inequalities and monotonicity in reciprocal
spectral variables. Octonionic products and Pfaffians realize the
same gradient-pairing table with Laplacian coefficients $32$ and $24$.
For their earlier harmonic-morphism and regular-minimality results, see
\cite[Section~4.5.2]{RiedlerThesis25} and
\cite[Corollary~5.2 and Remark~5.5]{HoppeTkachev19}.

\subsection{Complex determinants}
\label{subsec:intro-determinants}

Real parts of complex determinants have this gradient table with the
smaller Laplacian coefficient $16$. The preceding multiplier fails in
orders at least three. A product of proper minor norms supplies a
different certificate in every order, including across nonisolated
singularities. Here $E_j$ is the sum of squared absolute values of all
$j\times j$ minors and $F_m=\operatorname{Re}\det_{\mathbb C}Z$.

\begin{maintheorem}[Real parts of complex determinants]
\label{main:complex-determinants}
For every $m\ge1$, the boundary
$T_m=\partial\llbracket\{F_m<0\}\rrbracket$ is locally area
minimizing in $M_m(\mathbb C)\cong\R^{2m^2}$. For $m\ge2$,
\[
 G_m=F_m\left(\prod_{j=1}^{m-1}E_j\right)^{1/2}
\]
gives a subcalibration smooth off the rank-at-most-$m-2$ locus,
which has real codimension eight. For $m\ge2$, the cone is strictly area minimizing,
and
\begin{equation}\label{eq:main-determinant-stability}
 Q_{C_m}(u)\ge\frac{(m^2-m-1)^2}{4}\int_{C_m}r^{-2}u^2,
 \qquad u\in C_c^\infty((C_m)_{\reg}\setminus\{0\}).
\end{equation}
\end{maintheorem}
Regular minimality was proved in
\cite[Corollary~5.2 and Example~5.4]{HoppeTkachev19}.
For $m\ge3$, the smallest algebra has exactly $m+1$ independent
generators of maximal degree $2m$. The proof of global positivity
is uniform in $m$; for $m=3$ its potential is $F\sqrt{RS}$,
linking the spectral construction to the cubic one.

\Cref{sec:preliminaries,sec:laplacian-calculus,sec:algebraic-jacobi}
develop singular-set comparison, finite algebraic calculus and the
multiplier--Jacobi identity; \cref{sec:quartic-application} develops
pullback transfer. Structural restrictions in
\cref{sec:low-rank,sec:gradient-square-closure,sec:cubic-algebraic-theory}
support the quartic criteria and cubic minimizing classification in
\cref{sec:two-quadratic,sec:cubic-potentials,sec:nonradial-examples}.
\Cref{sec:spectral-positivity,sec:norm-products-pfaffians,sec:complex-determinants}
extend the method to arbitrary degree and unbounded generator number.

\begin{remark}[Smooth minimizing hypersurfaces]
These new area-minimizing algebraic hypercones also give smooth
minimizing hypersurfaces. Wang's extension
\cite[Theorem~1.1]{WangSmoothing24} of the Hardt--Simon theorem
\cite{HardtSimon85} provides, on each side of such a cone, a properly
embedded smooth area-minimizing hypersurface having that cone as its
tangent cone at infinity. This applies also to the examples with
nonisolated singularities.
\end{remark}
